
\begin{example}[Causal Delivery Mailbox\formalized{Examples/CausalMailbox.lean}{Examples.CausalMailbox}]\label{ex:causal-mailbox}\sloppy%
  Consider a distributed publish-subscribe system in which topic messages
  carry causal dependencies. A \emph{topic relay} engine forwards messages
  between nodes using the default FIFO mailbox, performing only signature
  validation. A \emph{local broker} engine delivers messages to subscribers
  on the same node, but must respect causal order: a message may only be
  delivered after all messages it depends on have been delivered. A FIFO
  mailbox cannot enforce this invariant because network reordering may
  deliver a message before its dependencies arrive.

  We model the system with two engine type indices:
  $\defequal{\EngIdx}{\{ \mathsf{relay}, \mathsf{broker} \}}$.
  Both share the message interface
  \begin{equation}\label{eq:topic-msg}
    \typedef{%
      \msgtype{\mathsf{relay}}%
    }{%
      \msgtype{\mathsf{broker}}%
    } \coloneqq
      \msg{TopicMsg}(\Address, \FinSet{\mathsf{Hash}}, \mathsf{Data},
      \mathsf{Sig}, \mathsf{Hash})\,,
  \end{equation}
  where the second component records the set of message hashes on which this
  message causally depends and the fifth component is the message's own hash.
  We write $\deps{v}$ for this dependency set
  and $\mhash{v}$ for the hash of message~$v$.

  The relay engine uses the default mailbox (\Cref{def:mailbox-engine}).
  The broker engine uses a custom mailbox with local state
  \begin{equation}\label{eq:causal-state}
    \typedef{%
      L_{\mailbox{\mathsf{broker}}}%
    }{\left(\!\!
    \begin{aligned}
      & \fdlv \colon \FinSet{\mathsf{Hash}},\\
      & \fpnd \colon \Map{\mathsf{Hash}}{\msg{TopicMsg}},\\
      & \frdy \colon \List{\msg{TopicMsg}}
    \end{aligned}
    \!\!\right)}\,,
  \end{equation}
  where $\fdlv$ records hashes of delivered messages,
  $\fpnd$ buffers messages whose dependencies have not yet been met,
  and $\frdy$ holds messages ready for delivery.

  The broker's mailbox behaviour consists of a single guarded action
  $\guardedactionval{\guard{c}}{\action{c}}$:
  \begin{itemize}
  \item The guard always matches:
    $\guard{c}(\append(v), k, q) \coloneqq \some{v}$.
  \item The action branches on dependency satisfaction.
    If $\deps{v} \subseteq q.\fdlv$, the message is ready
    and a cascade check releases any pending messages whose dependencies are
    now met:
    \begin{align*}
      \action{c}(v, \ldots) \coloneqq {}&\\
        \refeffect{Chain}\bigl(\;
          &\update{q[\fdlv \cup \{\mhash{v}\},\;
                    \frdy \pp [v]]},\\
          &\update{q[\fdlv \cup \{\mhash{v}\},\;
                    \fpnd \setminus S,\;
                    \frdy \pp [v] \pp S_{\mathsf{list}}]}\;\bigr)
    \end{align*}
    where $S = \{ w \in q.\fpnd \mid \deps{w}
    \subseteq q.\fdlv \cup \{\mhash{v}\}\}$.
    When $S = \emptyset$, the second effect coincides with the first;
    this case simplifies to $\chain{\update{\cdot}}{\noop}$.

    If $\deps{v} \not\subseteq q.\fdlv$, the message is
    buffered:
    \[
      \action{c}(v, \ldots) \coloneqq
        \update{q[\fpnd \cup \{\mhash{v} \mapsto v\}]}\,.
    \]
  \end{itemize}
  Because there is a single guarded action, the non-overlapping guard
  condition (\Cref{def:non-overlapping-guards}) is satisfied trivially. The
  use of $\chain{\cdot}{\cdot}$ demonstrates that mailbox behaviours may
  require effect sequencing to implement policies that go beyond simple
  message storage.
\end{example}
